\documentclass[10pt,a4paper]{amsart}
\usepackage[margin=19mm]{geometry}
\usepackage{amsmath,amssymb,mathtools}
\usepackage[scr=rsfs,cal=euler]{mathalfa}
\usepackage{xfrac}
\usepackage[unicode,hidelinks,bookmarks=false]{hyperref}
\newcommand{\ie}{i.e.\ }
\newcommand{\cf}{cf.\ }
\newcommand{\resp}{respectively\ }
\newcommand{\Subsection}{\subsection}
\providecommand{\nobreakdash}{\nobreak\hbox{-}}
\setlength{\emergencystretch}{2em}
\multlinegap0pt
\usepackage{amssymb}
\usepackage{algorithm}\usepackage{algpseudocode}
\usepackage{xcolor}
\newtheorem{lemma}{Lemma}
\title{An Inexact Augmented Lagrangian Method for $(L_0,L_1)$-Smooth Convex Optimization}
\author{Aleksandr Vyguzov, Fedor Stonyakin}
\date{Source: arXiv:2609.00261v1 (CC BY 4.0)}
\begin{document}
\maketitle

\noindent\emph{Source and licence.} An Inexact Augmented Lagrangian Method for $(L_0,L_1)$-Smooth Convex Optimization. Version arXiv:2609.00261v1 (CC BY 4.0), distributed by arXiv under the Creative Commons Attribution 4.0 International licence, http://creativecommons.org/licenses/by/4.0/, as recorded in the arXiv licence metadata for this exact version and shown on the abstract page https://arxiv.org/abs/2609.00261v1. Full licence notice: \texttt{licenses/arxiv-2609.00261-CC-BY-4.0.txt}.\par\smallskip

\noindent\textit{Editorial scope.} Counted unit: the article's lemma lem:fun\_convexity (source0.tex lines 627--633). The article's complete original proof of it is delivered with it (source0.tex lines 1581--1636, one \texttt{proof} environment of 56 lines, which the article places away from the statement). No mathematics is written, repaired or re-derived: the statement and the proof are the article's own lines, in the article's own notation, and no retained line is substituted or re-flowed.  Retained uncounted as the source-local context the proof needs: lines 465--473.  Nothing was omitted from any retained span, so the package carries no omission record.  No retained line contains a citation, so the wrapper carries no bibliography and no bibliography entry.  No retained line contains a figure environment, an image-inclusion command or drawing code, so no image is delivered and the package declares no asset dependency.  Editorial formatting only, in addition: an AMS \texttt{amsart} preamble that restates the article's own declarations and macros used by the retained lines, namely the environments \texttt{lemma} on separate counters, together with the standard packages those lines need.  The retained environments print in this document as Lemma 1 in order of appearance, whereas the article prints the same items with the numbering of its own full document.  The complete native source of the article as distributed in the arXiv e-print for this exact version is bundled unchanged as \texttt{sources/209011/source0.tex}.
		\begin{equation}\label{l_0def}
			L_0^{\mathcal{L}}
			=
			L_0^f
			+
			\beta\|A^\top A\|
			+
			L_1^f B_S,
		\end{equation}

	\begin{lemma}\label{lem:fun_convexity}
		Let
		\[
		g(x)=xL_0^{\mathcal L}(x),
		\]
		where $L_0^{\mathcal L}$ is defined in \eqref{l_0def}. Then $g$ is a convex and monotonically increasing function on $\mathbb{R}_+$.
	\end{lemma}

	\begin{proof}[Proof of Lemma~\ref{lem:fun_convexity}]
		We consider $g(x) = x L_0^{\mathcal{L}}(x)$, where 
		$$
		L_0^{\mathcal{L}}
		=
		L_0^f
		+
		\beta\|A^\top A\|
		+
		L_1^f B_S,
		$$
		with
		$B_S=\|A\|\bigl(\|y^*\|+\beta R_S\bigr)$
		and
		$R_S=\sqrt{\frac{2(\mathcal{L}(x_s)-f^\star)}{\beta}+\frac{\|y^*\|^2}{\beta^2}}+\frac{\|y^*\|}{\beta}.$
		
		Substituting the definition of $L_0^{\mathcal{L}}$ and denoting $\Delta = \mathcal{L}(x_s)-f^\star$ we get:
		
		$$
		g(x) = x \left( \left( L_0^f + 2 L_1^f \| A \| \| y^* \| \right) + \| A^\top A\| x + L_1^f \| A \| x \left(\sqrt{\frac{2 \Delta}{x} + \frac{\| y^* \|^2}{x^2} } + \frac{\| y^* \|}{x} \right) \right)
		$$
		
		Denoting $a_0 = L_0^f + 2 L_1^f \| A \| \| y^* \|, a_1 = \| A^\top A\|, a_2 = L_1^f \| A \|, a_3 = \| y^* \|^2, a_4 = 2 \Delta$ (notice that $a_i \geq 0$) we get
		\begin{equation}\label{lem:fun_convexity:eq_a}
			g(x) = a_1 x^2 + a_0 x + a_2 x \sqrt{a_4 x + a_3}
		\end{equation}
		
		By the definitions of the constants, we have $a_i \geq 0$ for all $i \in \{0, 1, 2, 3, 4\}$. 
		If $a_3 = a_4 = 0$, the function reduces to the quadratic polynomial $g(x) = a_1 x^2 + a_0 x$, which is trivially convex and monotonically increasing on $\mathbb{R}_{+}$. 
		Otherwise, $a_4 x + a_3 > 0$ for all $x > 0$. To establish the convexity and monotonicity of $g(x)$ on $\mathbb{R}_{+}$, we analyze its first and second derivatives on $\mathbb{R}_{++}$.
		
		Let $h(x) = x \sqrt{a_4 x + a_3}$. The first derivative of $h(x)$ is given by:
		\begin{equation*}
			h'(x) = \sqrt{a_4 x + a_3} + \frac{a_4 x}{2\sqrt{a_4 x + a_3}} = \frac{3a_4 x + 2a_3}{2\sqrt{a_4 x + a_3}}.
		\end{equation*}
		Since $a_3, a_4 \geq 0$ and $x > 0$, it follows that $h'(x) \geq 0$. 
		The first derivative of $g(x)$ is then:
		\begin{equation*}
			g'(x) = 2a_1 x + a_0 + a_2 h'(x).
		\end{equation*}
		As $a_0, a_1, a_2 \geq 0$, we conclude that $g'(x) \geq 0$ for all $x \in \mathbb{R}_{++}$. Since $g(x)$ is continuous on $\mathbb{R}_{+}$, this implies that $g(x)$ is monotonically increasing on $\mathbb{R}_{+}$.
		
		Next, we compute the second derivative of $h(x)$:
		\begin{equation*}
			h''(x) = \frac{a_4}{2\sqrt{a_4 x + a_3}} + \frac{a_4}{2\sqrt{a_4 x + a_3}} - \frac{a_4^2 x}{4(a_4 x + a_3)^{3/2}} = \frac{a_4}{(a_4 x + a_3)^{3/2}} \left( \frac{3a_4 x}{4} + a_3 \right).
		\end{equation*}
		Given that $a_3, a_4 \geq 0$ and $x > 0$, we have $h''(x) \geq 0$. 
		The second derivative of $g(x)$ is:
		\begin{equation*}
			g''(x) = 2a_1 + a_2 h''(x).
		\end{equation*}
		Since $a_1, a_2 \geq 0$ and $h''(x) \geq 0$, it follows that $g''(x) \geq 0$ for all $x \in \mathbb{R}_{++}$. 
		Because $g(x)$ is continuous on $\mathbb{R}_{+}$ and its second derivative is non-negative on $\mathbb{R}_{++}$, we conclude that $g(x)$ is convex on $\mathbb{R}_{+}$.
		
		This completes the proof.
	\end{proof}

\end{document}
